\section{Quais são os principais concorrentes destas empresas no que se refere ao item 3?}

No mercado de tecnologia, a NVIDIA enfrenta concorrência de várias empresas de grande porte, que disputam a liderança em setores como GPUs, IA, simulação e automação. Abaixo, estão listados os principais concorrentes em cada um dos produtos e serviços mencionados \cite{nvidia_competitors_ideiaveloz}:

\begin{itemize}
    \item \textbf{GeForce GPUs}: A linha GeForce da NVIDIA é amplamente reconhecida no mercado de \textit{gaming}, oferecendo desempenho gráfico avançado. No entanto, a \textbf{AMD} se destaca como sua principal concorrente com a linha de GPUs \textbf{Radeon}, que também é voltada para jogadores e entusiastas. Além disso, a \textbf{Intel} entrou recentemente no mercado de GPUs com sua linha \textbf{Intel Arc}, voltada para jogos e multimídia, aumentando a competitividade neste segmento.

    \item \textbf{Quadro e RTX para Profissionais}: No setor de GPUs profissionais, onde a NVIDIA lidera com as linhas Quadro e RTX, a concorrência é marcada pela \textbf{AMD} com sua linha \textbf{Radeon Pro}, que atende a demandas de design gráfico, modelagem 3D e simulação. A \textbf{Intel} também compete neste mercado com soluções como a linha Xeon, voltada para computação intensiva e visualização profissional.

    \item \textbf{Plataformas de IA e Computação em Nuvem}: A NVIDIA domina o mercado de IA com suas GPUs especializadas, como as séries A100 e H100, usadas em data centers para aprendizado profundo e análise de dados. Contudo, a \textbf{Google} oferece uma alternativa competitiva com o \textbf{Tensor Processing Unit (TPU)}, otimizado para cargas de trabalho de IA. A \textbf{Amazon Web Services (AWS)} e a \textbf{Microsoft Azure} também oferecem soluções para IA em nuvem, integrando opções de hardware de alto desempenho.

    \item \textbf{Omniverse}: A plataforma Omniverse da NVIDIA, voltada para simulação 3D e colaboração virtual, enfrenta concorrência de ferramentas populares como o \textbf{Unity} e o \textbf{Unreal Engine}, amplamente usados para desenvolvimento de mundos virtuais e simulação interativa. A \textbf{Autodesk}, com produtos como o \textit{Maya} e o \textit{3ds Max}, também oferece alternativas robustas para modelagem e simulação digital.

    \item \textbf{NVIDIA DRIVE para Veículos Autônomos}: A plataforma NVIDIA DRIVE integra hardware e software para veículos autônomos e sistemas ADAS. Neste segmento, a \textbf{Intel Mobileye} é uma das concorrentes principais, oferecendo soluções de visão computacional e sensores. Além disso, a \textbf{Waymo} (subsidiária da Alphabet) e a \textbf{Qualcomm} têm desenvolvido plataformas avançadas de software e hardware para condução autônoma e assistência ao motorista.
\end{itemize}
